\documentclass[10pt,a4paper]{amsart}
\usepackage[margin=19mm]{geometry}
\usepackage{amsmath,amssymb,mathtools}
\usepackage[scr=rsfs,cal=euler]{mathalfa}
\usepackage{xfrac}
\usepackage[unicode,hidelinks,bookmarks=false]{hyperref}
\newcommand{\ie}{i.e.\ }
\newcommand{\cf}{cf.\ }
\newcommand{\resp}{respectively\ }
\newcommand{\Subsection}{\subsection}
\providecommand{\nobreakdash}{\nobreak\hbox{-}}
\setlength{\emergencystretch}{2em}
\multlinegap0pt
\usepackage[shortlabels]{enumitem}
\usepackage{amssymb}
\usepackage{xcolor}
\usepackage{tcolorbox}
\newtheorem{Rq}{Remark}
\newtheorem{claim}{Claim}
\newcommand{\deffont}[1]{\textbf{#1}}
\title{Isomorphisms and automorphisms of graph products of groups}
\author{Amandine Escalier, Camille Horbez}
\date{Source: arXiv:2607.22303v2 (CC BY 4.0)}
\begin{document}
\maketitle

\noindent\emph{Source and licence.} Isomorphisms and automorphisms of graph products of groups. Version arXiv:2607.22303v2 (CC BY 4.0), distributed by arXiv under the Creative Commons Attribution 4.0 International licence, http://creativecommons.org/licenses/by/4.0/, as recorded in the arXiv licence metadata for this exact version and shown on the abstract page https://arxiv.org/abs/2607.22303v2. Full licence notice: \texttt{licenses/arxiv-2607.22303-CC-BY-4.0.txt}.\par\smallskip

\noindent\textit{Editorial scope.} Counted unit: the article's claim claim:step-1 (source0.tex lines 1101--1103). The article's complete original proof of it is delivered with it (source0.tex lines 1106--1165, one \texttt{proof} environment of 60 lines). No mathematics is written, repaired or re-derived: the statement and the proof are the article's own lines, in the article's own notation, and no retained line is substituted or re-flowed.  Retained uncounted as the source-local context the proof needs: lines 362--364.  Nothing was omitted from any retained span, so the package carries no omission record.  No retained line contains a citation, so the wrapper carries no bibliography and no bibliography entry.  No retained line contains a figure environment, an image-inclusion command or drawing code, so no image is delivered and the package declares no asset dependency.  Editorial formatting only, in addition: an AMS \texttt{amsart} preamble that restates the article's own declarations and macros used by the retained lines, namely the environments \texttt{Rq}, \texttt{claim} on separate counters, together with the standard packages those lines need.  The retained environments print in this document as Rq 1, Claim 1 in order of appearance, whereas the article prints the same items with the numbering of its own full document.  The complete native source of the article as distributed in the arXiv e-print for this exact version is bundled unchanged as \texttt{sources/205930/source0.tex}.
\begin{Rq}\label{Fact:Normalizer-support}
    If $N\subseteq G$ normalizes $A$, then it normalizes the parabolic support $R$ of $A$. In particular, we have $N_G(A)=N_R(A)\times R^\perp$.
\end{Rq}

\begin{claim}\label{claim:step-1}
The images $\varphi_i(S_i)$ and $\varphi_i(R_i)$ are $\overline{\Lambda}$-parabolic, whence $\Lambda$-parabolic.
\end{claim}

\begin{proof}[Proof of the claim]~

\textbf{Introducing tight subgroups.} The key point of this proof is the following notion which is invariant under isomorphisms: say that a subgroup $A\subseteq C_P^e$ is \deffont{tight} if for every subgroup $A'\subseteq C_P^e$ that properly contains $A$, one has $\langle C_P^e,N_G(A')\rangle\subsetneq \langle C_P^e,N_G(A)\rangle$. We define tight subgroups of $C_Q^e$ in the same way. Notice that:
\begin{enumerate}[label={(\alph*)}]
\item \label{item:tight1} If $A\subseteq C_P^e$ is a subgroup, and $B$ is the $\overline{\Gamma}$-parabolic support of $A$, Remark~\ref{Fact:Normalizer-support} gives $N_G(A)=N_B(A)\times B^\perp$, while $N_G(B)=B\times B^\perp$. As $B\subseteq C_P^e$, it follows that \begin{equation*}
    \left\langle C^e_P,N_G(A)\right\rangle
    =\left\langle C_P^e, B^\perp\right\rangle 
    \quad \text{and} \quad
    \left\langle C^e_P,N_G(B)\right\rangle
    =\left\langle C_P^e, B^\perp\right\rangle.
\end{equation*}
\item \label{item:tight2} Notice additionally that as the $\overline{\Gamma}$-type of $C_P^e$ is a clique, we have $C_P^e\subseteq B\times B^\perp$, so in fact $\langle C_P^e, B^\perp\rangle=B\times B^\perp=N_G(B)$.
\item \label{item:tight-parab} By point \ref{item:tight1}, tight subgroups of $C^e_P$ are necessarily $\overline{\Gamma}$-parabolic.
\item \label{item:tight-CNS} By points \ref{item:tight1} and~\ref{item:tight2}, a parabolic subgroup $A$ is tight if and only if for every \emph{$\mathit{\overline{\Gamma}}$-parabolic} subgroup $A'\subseteq C^e_P$ that properly contains $A$, one has $N_G(A')\subsetneq N_G(A)$.
\end{enumerate} 
\medskip

\textbf{Showing that {$S_i$ is tight and} $\varphi_i(S_i)$ is $\overline{\Lambda}$-parabolic.} We now prove that $S_i$ is tight, using~\ref{item:tight-CNS}. Since tightness is preserved by isomorphism, this also implies that $\varphi_i(S_i)$ is tight, whence $\overline{\Lambda}$-parabolic by \ref{item:tight-parab}.

So let $U\subseteq C_P^e$ be a parabolic subgroup that properly contains $S_i$. By definition of $S_i$ and since $S_i \subsetneq U$, there exists a vertex group $D\subseteq U$ such that $N_G(P_i)\not\subseteq N_G(D)$. Since $N_G(P_i)=N_G(S_i)$, this rewrites as $N_G(S_i)\not\subseteq N_G(D)$, or in other words 
\begin{equation}\label{eq:1}
    N_G(S_i)\cap N_G(D)\subsetneq N_G(S_i).
\end{equation}
Let $\Delta^*_U$ (resp.\ $\Delta^*_{S_i}$, resp.\ $\Delta^*_D$) be the full subgraph of $\overline{\Gamma}$ whose vertex set consists of all vertices that are in the type of $U^\perp$ (resp.\ $S_i^\perp$, resp.\ $D^\perp$) but not in the type of $C_P^e$. Then,
\begin{equation*}
    N_G(U)=\left\langle{C_P^e, G_{\Delta_U^*}} \right\rangle, \quad 
    N_G(S_i)=\left\langle{C_P^e, G_{\Delta_{S_i}^*}} \right\rangle, \quad 
    N_G(D)=\left\langle{C_P^e, G_{\Delta_D^*}} \right\rangle.
\end{equation*} 
As $D\subseteq U$ and $S_i\subseteq U$, we have $\Delta_U^*\subseteq \Delta_D^*$ and $\Delta_U^*\subseteq \Delta_{S_i}^*$. Therefore 
\begin{equation}\label{eq:2}
    N_G(U)\subseteq N_G(S_i)\cap N_G(D).
\end{equation}
Combining Equations~\eqref{eq:1} and~\eqref{eq:2} shows that $N_G(U)\subsetneq N_G(S_i)$, so by \ref{item:tight-CNS} $S_i$ is tight.

\medskip

\textbf{Showing that $\varphi_i(R_i)$ is $\overline{\Lambda}$-parabolic.}
Let us first prove that every vertex group $B\subseteq R_i$ is contained in a tight parabolic subgroup $T_B\subseteq R_i$.

So let $T_B:=B^+\cap S_i$, in other words $T_B$ is the maximal parabolic subgroup of $S_i$ with $N_G(B)=N_G(T_B)$. Remark that since $B\subseteq S_i$ and $N_G(S_i)=N_G(P_i)$, we have 
\begin{equation}
    \label{eq:normalisateur-B-NGPi}
    N_G(P_i)\subseteq N_G(B).
\end{equation}
\begin{itemize}
    \item As $\overline{\Gamma}$ is strongly reduced, the vertices of $\overline{\Gamma}$ representing the types of $P_i$ and $B$ are not equivalent. Equation~\eqref{eq:normalisateur-B-NGPi} thus implies that $N_G(P_i)\subsetneq N_G(B)$.
    Therefore $T_B$ does not contain $P_i$, in particular $T_B$ is a proper $\overline{\Gamma}$-parabolic subgroup of $S_i$.
    Recalling that $S_i=P_i\times R_i$ shows that $T_B\subseteq R_i$. 
    \item We claim that $T_B$ is tight. If not, then by the above point \ref{item:tight-CNS}, $T_B$ is properly contained in a parabolic subgroup $R\subseteq C_P^e$ such that $N_G(R)=N_G(T_B)$. By maximality of $T_B$ inside $S_i$, necessarily $R$ is not contained in $S_i$. Let $\Delta_R^*$ (resp.\ $\Delta_{S_i}^*$) be the full subgraph of $\overline{\Gamma}$ whose vertex set consists of all vertices that are in the type of $R^\perp$ (resp.\ $S_i^\perp$) but not in the type of $C_P^e$. Then 
\begin{equation}\label{eq:3}
    N_G(S_i)=\left\langle { C_P^e, G_{\Delta_{S_i}^*} }\right\rangle, 
    \quad N_G\left( \left\langle{S_i,R}\right\rangle \right)=\left\langle {C_P^e, G_{\Delta_R^*\cap \Delta_{S_i}^*} } \right\rangle.
\end{equation}
On the other hand, Equation~\eqref{eq:normalisateur-B-NGPi} rewrites as $N_G(S_i)\subseteq N_G(R)$, which shows that $\Delta_{S_i}^*\subseteq \Delta_R^*$. Equation~\eqref{eq:3} therefore gives $N_G(S_i)=N_G(\langle S_i,R\rangle)$, contradicting the fact that $S_i$ is tight. So $T_B$ is tight.
\end{itemize}
Hence, every vertex group $B\subseteq R_i$ is contained in a tight parabolic subgroup $T_B\subseteq R_i$. 

This implies that $R_i$ is generated by tight subgroups of $C_P^e$, and therefore $\varphi_i(R_i)$ is generated by tight subgroups of $C_Q^e$. This and \ref{item:tight-parab}, together with the fact that the type of $C_Q^e$ is a clique, shows that $\varphi_i(R_i)$ is $\overline{\Lambda}$-parabolic.
\end{proof}

\end{document}
